\noindent\underline{Theorem 4}\quad
If $(X_i)$ iid $P_{\theta_0}$, under regularity conditions,
\[
\sqrt{n}(\hat{\theta}_n-\theta_0)\xrightarrow{d}N(0,I^{-1}(\theta_0)).
\]

\noindent\underline{Proof}\quad
Define $E_n=\{\hat{\theta}_n\in K\}$, for a ``regular'' neighbourhood of $\theta_0$.
By consistency $\mathbb{P}(E_n)\to1$ as $n\to\infty$.
Working on $E_n$, a Taylor series gives
\[
\sqrt{n}(\hat{\theta}_n-\theta_0)
=-(\overline{A}_n)^{-1}
\underbrace{\sqrt{n}\,\nabla_\theta\overline{\ell}_n(\theta_0)}
_{\xrightarrow{d}N(0,I(\theta_0))\text{ by a CLT}}
=:Z_n,
\]
where
\[
(\overline{A}_n)_{j,k}
=\frac1n\sum_{i=1}^n
\underbrace{\frac{\partial^2}{\partial\theta_j\partial\theta_k}
\log f(X_i,\theta^{(j)})}_{=:g(X_i,\theta^{(j)})},
\qquad\theta^{(j)}\in[\hat{\theta}_n,\theta_0].
\]
\noindent\underline{Claim}\quad
$\overline{A}_n\xrightarrow{P}-I(\theta_0)$.

Several applications of Slutsky's Lemma and CMT show that
$Z_n\xrightarrow{d}N(0,I^{-1}(\theta_0))$.

Why can we work on $E_n$? On $E_n$,
$\sqrt{n}(\hat{\theta}_n-\theta_0)=Z_n$. Hence,
\[
\sqrt{n}(\hat{\theta}_n-\theta_0)
=\underbrace{Z_n}_{\xrightarrow{d}N(0,I^{-1}(\theta_0))}
+\underbrace{\bigl(\sqrt{n}(\hat{\theta}_n-\theta_0)-Z_n\bigr)}
_{\xrightarrow{P}0},
\]
as
\[
\mathbb{P}\bigl(\lvert\sqrt{n}(\hat{\theta}_n-\theta_0)-Z_n\rvert<\varepsilon\bigr)
\geq\mathbb{P}(E_n)\xrightarrow{n\to\infty}1.
\]
By Slutsky,
$\sqrt{n}(\hat{\theta}_n-\theta_0)\xrightarrow{d}N(0,I^{-1}(\theta_0))$.

\medskip
\noindent\underline{Proof of Claim}\quad
We need to show, $\forall j,k$,
\[
\lvert(\overline{A}_n)_{j,k}-(-I(\theta_0))_{j,k}\rvert\xrightarrow{P}0.
\]
Let $g(\theta)=\mathbb{E}_{\theta_0}[g(X,\theta)]$.
Recall that $I(\theta_0)=-g(\theta_0)$ (under reg conditions).
Note
\[
(\overline{A}_n)_{j,k}+I(\theta_0)_{j,k}
=\frac1n\sum_{i=1}^n g(X_i,\theta^{(j)})-g(\theta^{(j)})
+g(\theta^{(j)})-g(\theta_0).
\]
Then,
\[
\begin{aligned}
\lvert(\overline{A}_n)_{j,k}+I(\theta_0)_{j,k}\rvert
&\leq\underbrace{\sup_{\theta\in K}
\left\lvert\frac1n\sum_{i=1}^n g(X_i,\theta)-g(\theta)\right\rvert}_{\text{(I)}}
+\underbrace{\lvert g(\theta^{(j)})-g(\theta_0)\rvert}_{\text{(II)}}.
\end{aligned}
\]
By a ULLN, $\text{(I)}\xrightarrow{\mathrm{a.s.}}0$, using
\[
\theta\longmapsto g(x,\theta)
=\frac{\partial^2}{\partial\theta_j\partial\theta_k}\log f(x,\theta)
\quad\text{is continuous},
\]
$K$ is compact, and
$\mathbb{E}_{\theta_0}\sup_{\theta\in K}|g(X,\theta)|<\infty$.
Also, we know $\hat{\theta}_n\xrightarrow{P}\theta_0$ by consistency,
thus $\theta^{(j)}\xrightarrow{P}\theta_0$, and an application of the
dominated convergence theorem shows that $g$ is continuous in $K$.
Hence, by the CMT,
$\text{(II)}\xrightarrow{P}0$.
Therefore,
\[
\lvert(\overline{A}_n)_{j,k}+I(\theta_0)_{j,k}\rvert\xrightarrow{P}0.
\qquad\square
\]

\noindent (1) The continuous derivatives of
$\theta\mapsto\log f(x,\theta)$ are \underline{not} necessary for asymptotic normality.
Le Cam theory relaxes this for models with a single weak derivative.

\medskip
\noindent\underline{Ex}\quad Laplace model
\[
f(x,\theta)=\frac12\exp\{-|x-\theta|\},\qquad x\in\mathbb{R}.
\]
\noindent (2) One weak derivative is necessary.

\medskip
\noindent\underline{Ex}\quad $X\sim\operatorname{Unif}([0,\theta])$, $\theta>0$.

The MLE is $\hat{\theta}_n=\max_{1\leq i\leq n}X_i$, and
\[
f(x,\theta)=\frac1\theta\,\mathbf{1}_{\{x\leq\theta\}}.
\]
\begin{center}
\setlength{\unitlength}{0.6mm}
\begin{picture}(65,22)
\put(0,3){\line(1,0){60}}
\qbezier(3,3)(14,3)(22,4)
\put(22,4){\line(0,1){16}}
\qbezier(22,20)(24,9)(36,6)
\qbezier(36,6)(48,3)(60,2)
\end{picture}
\end{center}
The likelihood is discontinuous at $\theta=x$ and
$\sqrt{n}(\hat{\theta}_n-\theta_0)$ has a non-normal limit (Ex sheet).

\medskip
\noindent (3) It is also necessary that $\theta_0\in\operatorname{int}(\Theta)$.

\medskip
\noindent\underline{Ex}\quad $X\sim N(\theta,1)$, where $\Theta=[0,\infty)$.

Here the MLE $\hat{\theta}_n=\max(0,\overline{X}_n)$.
Under $P_{\theta_0}$ where $\theta_0=0$ (on the boundary of $[0,\infty)$),
we can show that $\sqrt{n}(\hat{\theta}_n-\theta_0)$ has a non-normal limit (Ex sheet).

\medskip
\noindent (4) Cramer Rao lower bound says that $I^{-1}(\theta_0)$ is
best variance for $\sqrt{n}(\hat{\theta}_n-\theta_0)$ for unbiased
$\hat{\theta}_n$, for any finite $n$.

\medskip
\noindent\underline{Le Cam theory}\quad
In models with the ``local asymptotic normality'', any estimator
$\tilde{\theta}_n$ which is ``regular'' has asymptotic variance
$\geq I^{-1}(\theta_0)$.
Non-regular estimators can have better asymptotic variance for certain values of $\theta_0$.

\medskip
\noindent\underline{Ex}\quad Hodge's estimator

Let $\hat{\theta}_n$ be the MLE in a regular model s.t.
\[
\sqrt{n}(\hat{\theta}_n-\theta_0)\xrightarrow{d}N(0,I^{-1}(\theta_0))
\qquad\forall\theta_0\in\operatorname{int}(\Theta).
\]
Define
\[
\tilde{\theta}_n=
\begin{cases}
0&\text{if }|\hat{\theta}_n|<n^{-1/4},\\
\hat{\theta}_n&\text{o.w.}
\end{cases}
\]
If $\theta_0\neq0$, under $P_{\theta_0}$,
\[
\begin{aligned}
\mathbb{P}(\tilde{\theta}_n\neq\hat{\theta}_n)
&=\mathbb{P}\bigl(\underbrace{|\hat{\theta}_n|}
_{|\theta_0|-|\hat{\theta}_n-\theta_0|\leq}<n^{-1/4}\bigr)\\
&\leq\mathbb{P}\bigl(|\theta_0|-|\hat{\theta}_n-\theta_0|<n^{-1/4}\bigr)\\
&=\mathbb{P}\bigl(|\hat{\theta}_n-\theta_0|>|\theta_0|-n^{-1/4}\bigr)
\xrightarrow{n\to\infty}0
\qquad\text{as }\hat{\theta}_n\text{ consistent}.
\end{aligned}
\]
Therefore, under $P_{\theta_0}$,
$\sqrt{n}(\tilde{\theta}_n-\theta_0)\xrightarrow{d}N(0,I^{-1}(\theta_0))$.
However, if $\theta_0=0$,
\[
\begin{aligned}
\mathbb{P}(\hat{\theta}_n=\tilde{\theta}_n)
&=\mathbb{P}(|\hat{\theta}_n|\geq n^{-1/4})\\
&=\mathbb{P}\bigl(\sqrt{n}\,|\hat{\theta}_n-\theta_0|\geq n^{1/4}\bigr).
\end{aligned}
\]
Recall that by Prop 1 in lecture 4,
$\sqrt{n}(\hat{\theta}_n-\theta_0)=O_{\mathbb{P}}(1)$, which implies
the RHS $\xrightarrow{n\to\infty}0$.
In other words,
\[
\mathbb{P}(\tilde{\theta}_n=0)\xrightarrow{n\to\infty}1
\quad\text{or}\quad
\sqrt{n}(\tilde{\theta}_n-\theta_0)\xrightarrow{d}0.
\]
Hence the asymptotic variance is $0<I^{-1}(\theta_0)$ in general.

\medskip
\noindent\underline{Mean squared error}

\begin{center}
\setlength{\unitlength}{0.9mm}
\begin{picture}(140,60)
\put(36,9){\vector(1,0){100}}
\put(84,8){\vector(0,1){48}}
\put(84,4){\makebox(0,0){$0$}}
\put(135,4){\makebox(0,0){$\theta$}}
\put(65,40){\makebox(0,0){$n^{-1}$}}
\put(36,27){\line(1,0){100}}
\put(8,32){\makebox(0,0){$\mathbb{E}_\theta(|\hat{\theta}_n-\theta|^2)$}}
\qbezier(36,27)(54,26)(59,31)
\qbezier(59,31)(65,37)(68,47)
\qbezier(68,47)(73,54)(76,25)
\qbezier(76,25)(80,9)(84,9)
\qbezier(84,9)(89,9)(94,35)
\qbezier(94,35)(99,64)(106,41)
\qbezier(106,41)(112,27)(122,28)
\qbezier(122,28)(130,29)(136,27)
\put(119,47){\makebox(0,0){$\mathbb{E}_\theta[(\tilde{\theta}_n-\theta)^2]$}}
\end{picture}
\end{center}
By a ``minimax'' criterion, the Hodges estimator is suboptimal.

\begin{center}
\textbf{End of Lecture 7}
\end{center}
